\chapter{Il Web sotto il cofano}

Finora abbiamo guardato il Web dal lato del browser: HTML, CSS, JavaScript, e
infine \code{fetch} per parlare con un server. Prima di passare dall'altra parte
e scrivere noi il server, apriamo il cofano e guardiamo che cosa succede
\textbf{fra} il browser e il server: come comunicano due programmi su macchine
diverse, come si trova un server partendo da un nome, che cosa dice esattamente
HTTP e come lo protegge HTTPS.

Le basi (il ciclo richiesta--risposta, le URL, la struttura dei messaggi, i
codici di stato) le abbiamo viste nel capitolo 1: qui le riprendiamo e le
approfondiamo. L'obiettivo dichiarato della lezione è che, alla fine, la scheda
\textit{Rete} dei DevTools faccia un po' meno paura.

% ============================================================
\section{Le applicazioni di rete}

\dfn{Applicazione di rete}{
  Un'\term{applicazione di rete} è composta da \textbf{più programmi} che girano
  su sistemi diversi (gli \term{host}, o \textit{end system}) e comunicano fra
  loro attraverso la rete.}

Un'applicazione web, ad esempio, comprende due programmi distinti:

\begin{itemize}
  \item il \textbf{browser}, che gira sull'host dell'utente (desktop, portatile,
        tablet, smartphone\ldots);
  \item il \textbf{server web}, che gira sull'host del server.
\end{itemize}

\subsection{Reti di scala diversa}

Le reti si classificano per estensione geografica. Cambiano la scala e chi le
gestisce, ma grazie ai \textbf{protocolli} comuni reti di ogni scala riescono a
comunicare fra loro: Internet è, letteralmente, una \textbf{rete di reti}.

\begin{figure}[H]
\centering
\begin{tikzpicture}[
    chip/.style={rounded corners=3pt, fill=#1, text=white, font=\Lato\scriptsize\bfseries,
      minimum width=1.25cm, minimum height=0.55cm},
    riga/.style={rounded corners=3pt, draw=neutral!35!white, fill=neutral!4!white,
      line width=0.6pt, minimum width=15.6cm, minimum height=0.8cm},
    est/.style={font=\Lato\small\bfseries, anchor=west},
    des/.style={font=\Lato\small, anchor=west, text=neutral!25!black}]
  \foreach \i/\sigla/\col/\estensione/\descr in {
      0/PAN/orange-gradient-6/{pochi metri}/{i dispositivi personali (Bluetooth: telefono, cuffie, orologio)},
      1/LAN/accent!75!black/{un edificio, un campus}/{alta velocità, gestita localmente (la rete di casa o del dipartimento)},
      2/MAN/green-gradient-6/{una città, un'area metropolitana}/{collega più LAN},
      3/WAN/primary!75!black/{nazioni, continenti}/{gli operatori interconnettono le proprie reti}} {
    \node[riga, anchor=west] (r\i) at (0,-\i*1.0) {};
    \node[chip=\col] at (0.85,-\i*1.0) {\sigla};
    \node[est, text=\col] at (1.75,-\i*1.0) {\estensione};
    \node[des] at (6.9,-\i*1.0) {\descr};
  }
\end{tikzpicture}
\caption{Le reti per scala. Internet le collega tutte.}
\end{figure}

\subsection{Client-server e peer-to-peer}

L'\term{architettura} di un'applicazione di rete la decide lo sviluppatore, e
stabilisce come l'applicazione è distribuita sui vari host. I modelli
fondamentali sono due.

\begin{itemize}
  \item \textbf{Client-server}: i nodi sono \textbf{eterogenei}. C'è un host
        \textbf{sempre acceso}, il \term{server}, che fornisce servizi su
        richiesta di molti altri host, i \term{client}.
  \item \textbf{Peer-to-peer} (P2P): i nodi sono (idealmente)
        \textbf{omogenei}. L'applicazione sfrutta la comunicazione diretta fra
        coppie di host connessi in modo intermittente, i \term{peer}.
\end{itemize}

\begin{figure}[H]
\centering
\begin{tikzpicture}[
    srv/.style={wtnode, fill=primary!18!white, minimum width=1.5cm, minimum height=0.8cm,
      font=\Lato\scriptsize\bfseries},
    cli/.style={wtnodeplain, minimum width=1.2cm, font=\Lato\scriptsize},
    peer/.style={wtnodeacc, minimum width=1.0cm, font=\Lato\scriptsize},
    lnk/.style={{Stealth[length=4pt,width=4pt]}-{Stealth[length=4pt,width=4pt]},
      line width=0.7pt, draw=primary!65!black},
    plnk/.style={line width=0.7pt, draw=accent!60!black, dashed}]
  % --- client-server ---
  \node[srv] (s) at (0,0) {\icpu\ server};
  \foreach \a/\n in {90/1, 162/2, 234/3, 306/4, 18/5} {
    \node[cli] (c\n) at (\a:2.1) {client};
    \draw[lnk] (c\n) -- (s);
  }
  \node[font=\Lato\small\bfseries] at (0,-2.85) {Client-server};
  % --- peer-to-peer ---
  \begin{scope}[shift={(7.6,0)}]
    \foreach \a/\n in {90/1, 150/2, 210/3, 270/4, 330/5, 30/6} \node[peer] (p\n) at (\a:1.9) {peer};
    \foreach \x/\y in {1/3, 1/4, 2/5, 3/6, 4/6, 2/4, 5/1} \draw[plnk] (p\x) -- (p\y);
    \node[font=\Lato\small\bfseries] at (0,-2.85) {Peer-to-peer};
  \end{scope}
\end{tikzpicture}
\caption{A sinistra i client parlano solo con il server; a destra ogni peer può
         parlare direttamente con gli altri.}
\end{figure}

Nell'architettura client-server:

\begin{itemize}
  \item i client \textbf{non} comunicano direttamente fra loro;
  \item il server ha un indirizzo \textbf{fisso e noto} (un nome di host o un
        indirizzo IP), così un client può sempre contattarlo inviandogli una
        richiesta;
  \item per reggere tutte le richieste, spesso i server sono \textbf{molti}, ma
        il client di solito non se ne accorge e li percepisce come un unico
        server.
\end{itemize}

I server multipli possono essere raccolti in \term{data center} (enormi
quantità di server in un unico luogo, che vanno alimentati, mantenuti e ben
collegati), sparsi come \textbf{server distribuiti} in tutto il mondo (che vanno
interconnessi e coordinati), oppure organizzati in \textbf{data center
distribuiti}. È il caso di Google, che secondo i dati riportati a lezione conta
fra 30 e 50 data center nel mondo, che insieme gestiscono la ricerca ma anche
YouTube, Gmail e gli altri servizi.

Un'applicazione web è l'esempio tipico di client-server: c'è un server web
\textbf{sempre acceso e raggiungibile}, che riceve le richieste dei browser e
risponde inviando l'oggetto richiesto.

% ============================================================
\section{Processi che comunicano}

In un'applicazione di rete, a comunicare non sono propriamente gli host, ma i
\term{processi} che girano su di essi, magari su sistemi operativi diversi.

\begin{itemize}
  \item In un'applicazione web, il processo del browser del client scambia
        messaggi con il processo del server web.
  \item In un'applicazione P2P di \textit{file sharing}, un file viene trasferito
        dal processo di un peer al processo di un altro peer.
\end{itemize}

In ogni coppia di processi che comunicano c'è tipicamente un \textbf{processo
client}, quello che avvia la comunicazione, e un \textbf{processo server},
quello che attende di essere contattato. Nel Web il browser è il client e il
server web è il server; nel P2P il peer che scarica fa da client e quello che
invia fa da server. In un'applicazione P2P, quindi, un processo può
\textbf{cambiare ruolo} da uno scambio all'altro.

\subsection{Il socket}

\dfn{Socket}{
  Un \term{socket} è l'\textbf{interfaccia software} fra il livello applicativo e
  il livello di trasporto, attraverso cui un processo invia e riceve messaggi
  sulla rete. Nello stack TCP/IP è l'interfaccia fra l'applicazione (HTTP, DNS,
  SMTP\ldots) e il protocollo di trasporto (TCP, UDP o QUIC).}

\begin{figure}[H]
\centering
\begin{tikzpicture}[
    liv/.style={rounded corners=3pt, fill=#1, text=white, align=center,
      minimum width=3.0cm, minimum height=1.25cm, font=\Lato\scriptsize},
    fr/.style={-{Stealth[length=6pt,width=5pt]}, line width=1.1pt, draw=neutral!55!black}]
  \node[liv=accent!75!black] (a) at (0,0)    {\textbf{\small Applicazione}\\ HTTP, DNS, SMTP\ldots};
  \node[liv=orange-gradient-6] (s) at (3.7,0) {\textbf{\small Socket}\\ interfaccia software};
  \node[liv=green-gradient-6] (t) at (7.4,0)  {\textbf{\small Trasporto}\\ TCP / UDP / QUIC};
  \node[liv=primary!75!black] (n) at (11.1,0) {\textbf{\small Rete}\\ IP e instradamento};
  \draw[fr] (a) -- (s); \draw[fr] (s) -- (t); \draw[fr] (t) -- (n);
  \draw[wtbrace] ([yshift=4pt]a.north east) -- ([yshift=4pt]a.north west)
    node[midway, above=4pt, font=\Lato\scriptsize, align=center]{controllato dallo\\ sviluppatore};
  \draw[wtbrace] ([yshift=4pt]n.north east) -- ([yshift=4pt]t.north west)
    node[midway, above=4pt, font=\Lato\scriptsize, align=center]{controllato dal sistema operativo:\\
      lo sviluppatore sceglie solo il protocollo e pochi parametri};
\end{tikzpicture}
\caption{Il socket è la porta fra ciò che scriviamo noi e ciò che fa la rete.}
\end{figure}

Un'analogia classica è quella della \textbf{cassetta postale}. Quando un
processo vuole inviare un messaggio a un processo su un altro host, lo
imbuca nella propria cassetta (il socket). Il processo mittente dà per scontato
che, dall'altra parte della porta, ci sia un'infrastruttura di trasporto che
porterà il messaggio fino alla cassetta del destinatario. Arrivato all'host di
destinazione, il messaggio passa per la cassetta (il socket) del processo
ricevente.

\begin{figure}[H]
\centering
\begin{tikzpicture}[
    proc/.style={wtnode, minimum width=2.2cm, font=\Lato\scriptsize},
    cass/.style={rounded corners=2pt, draw=red-gradient-6, fill=red-gradient-1!60!white,
      line width=0.8pt, minimum width=1.0cm, minimum height=0.6cm, font=\Lato\tiny\bfseries},
    viaggio/.style={-{Stealth[length=6pt,width=5pt]}, line width=1pt, dashed,
      draw=red-gradient-6}]
  \node[proc] (p1) at (0,0) {processo\\ sull'host 1};
  \node[cass] (c1) at (2.4,0) {socket};
  \node[proc] (p2) at (11.2,0) {processo\\ sull'host 2};
  \node[cass] (c2) at (8.8,0) {socket};
  \draw[wtarrow] (p1) -- node[wtlab, above=1pt]{imbuca} (c1);
  \draw[wtarrow] (c2) -- node[wtlab, above=1pt]{ritira} (p2);
  \draw[viaggio] (c1.east) to[out=-25, in=205]
    node[wtlab, below=2pt]{infrastruttura di trasporto (la rete)} (c2.west);
  \draw[wtdash] (3.1,-1.25) -- (3.1,1.0) node[above, font=\Lato\tiny, text=neutral!40!black]{la «porta di casa»};
  \draw[wtdash] (8.1,-1.25) -- (8.1,1.0);
\end{tikzpicture}
\caption{Il socket come cassetta postale: il processo imbuca, la rete consegna.}
\end{figure}

Il socket è l'\textbf{API} fra l'applicazione e la rete. Lo sviluppatore
controlla tutto ciò che sta dal lato applicativo, ma ha poco controllo sul lato
del trasporto: può soltanto
\begin{enumerate}
  \item \textbf{scegliere il protocollo} di trasporto (TCP o UDP);
  \item forse \textbf{impostare qualche parametro} (dimensione massima dei
        buffer, dimensione massima dei segmenti\ldots).
\end{enumerate}
Una volta scelto il protocollo, l'applicazione si costruisce dando per scontati
i servizi che quel protocollo offre.

\subsection{Indirizzi IP e porte}

Per spedire una lettera serve un indirizzo. Su un host girano molte
applicazioni di rete contemporaneamente, quindi per raggiungere il processo
giusto servono \textbf{due} informazioni:

\begin{enumerate}
  \item l'\textbf{indirizzo dell'host}, per trovare la macchina giusta nella rete:
        è l'\term{indirizzo IP}, che lavora a livello di \textbf{rete}
        (instradamento);
  \item un \textbf{identificatore del processo}, per trovare il processo giusto
        dentro la macchina: è il \term{numero di porta}, che lavora a livello di
        \textbf{trasporto}.
\end{enumerate}

\begin{figure}[H]
\centering
\begin{tikzpicture}[
    host/.style={rounded corners=5pt, draw=primary!55!black, fill=primary!5!white,
      line width=0.8pt},
    proc/.style={wtnodeplain, minimum width=2.4cm, font=\Lato\scriptsize},
    porta/.style={rectangle, draw=red-gradient-6, fill=red-gradient-1!60!white, line width=0.7pt,
      font=\FiraCodeNL\scriptsize\bfseries, inner sep=2pt, minimum width=0.9cm},
    ip/.style={font=\FiraCodeNL\small\bfseries, text=primary!75!black}]
  % host client
  \draw[host] (0,-1.6) rectangle (3.6,1.4);
  \node[ip, anchor=south] at (1.8,1.45) {192.168.1.7};
  \node[proc, fill=accent!10!white] (b) at (1.5,0.4) {browser};
  \node[proc] at (1.5,-0.7) {client di posta};
  \node[porta] (pb) at (3.6,0.4) {52314};
  \draw[line width=0.6pt] (b) -- (pb);
  % host server
  \draw[host] (8.2,-1.6) rectangle (11.8,1.4);
  \node[ip, anchor=south] at (10.0,1.45) {93.184.216.34};
  \node[proc, fill=accent!10!white] (w) at (10.3,0.4) {server web};
  \node[proc] (m) at (10.3,-0.7) {server di posta};
  \node[porta] (p80) at (8.2,0.4) {80};
  \node[porta] (p25) at (8.2,-0.7) {25};
  \draw[line width=0.6pt] (w) -- (p80);
  \draw[line width=0.6pt] (m) -- (p25);
  % messaggio
  \draw[wtarrowacc] (pb) -- node[wtlab, above=2pt, align=center]
    {all'host \code{93.184.216.34},\\ al processo sulla porta \code{80}} (p80);
  \node[font=\Lato\scriptsize, text=neutral!35!black, align=center] at (5.9,-1.1)
    {l'IP trova la macchina,\\ la porta trova il processo};
\end{tikzpicture}
\caption{Indirizzo IP e porta. Anche il browser usa una porta, assegnata al
         momento dal sistema operativo (qui \code{52314}): serve a ricevere la
         risposta.}
\end{figure}

\begin{itemize}
  \item Un host è identificato da un indirizzo IP. Un indirizzo IPv4, come
        \code{192.168.1.10}, è un numero a \textbf{32 bit}; gli indirizzi IPv6,
        che stanno progressivamente sostituendo IPv4, sono a 128 bit.
  \item Un processo è identificato da un numero di porta (da 0 a 65535). Alle
        applicazioni più diffuse sono state assegnate per convenzione porte
        \textbf{note} (\textit{well-known}): la IANA ne pubblica l'elenco
        ufficiale.
\end{itemize}

\begin{center}
\renewcommand{\arraystretch}{1.3}
\begin{tabular}{@{}l@{\hspace{0.9cm}}l@{\hspace{1.6cm}}l@{\hspace{0.9cm}}l@{}}
  \textbf{Porta} & \textbf{Servizio} & \textbf{Porta} & \textbf{Servizio} \\
  \code{80}  & HTTP (Web)          & \code{25}  & SMTP (invio di posta) \\
  \code{443} & HTTPS (Web sicuro)  & \code{53}  & DNS \\
  \code{22}  & SSH                 & \code{21}  & FTP \\
\end{tabular}
\end{center}

È il motivo per cui nelle URL la porta si può omettere: se non la scriviamo, il
browser usa quella nota del protocollo (80 per \code{http}, 443 per
\code{https}).

% ============================================================
\section{I protocolli applicativi}

\dfn{Protocollo applicativo}{
  Un \term{protocollo di livello applicativo} definisce come i processi di
  un'applicazione di rete, in esecuzione su host diversi, si scambiano messaggi.
  In particolare stabilisce:
  \begin{enumerate}
    \item i \textbf{tipi} di messaggio scambiati (ad esempio richieste e risposte);
    \item la \textbf{sintassi} di ciascun tipo: quali campi contiene e come sono
          delimitati;
    \item la \textbf{semantica} dei campi, cioè il significato delle informazioni
          che contengono;
    \item le \textbf{regole} che stabiliscono quando e come un processo invia
          messaggi e risponde.
  \end{enumerate}}

In rete, e su Internet, si usano comunemente molti protocolli applicativi.
Alcune applicazioni si scambiano informazioni \textbf{sensibili} (credenziali,
dati personali\ldots): su una rete pubblica i messaggi vanno protetti, e per
questo molti protocolli hanno una \textbf{versione sicura}, evidenziata in verde.

\begin{center}
\renewcommand{\arraystretch}{1.35}
\begin{tabular}{@{}l@{\hspace{0.8cm}}>{\raggedright\arraybackslash}p{11.5cm}@{}}
  \textbf{Protocollo} & \textbf{A che cosa serve} \\
  DHCP & \textit{Dynamic Host Configuration Protocol}: assegna gli indirizzi IP agli host \\
  DNS  & \textit{Domain Name System}: traduce i nomi dei siti in indirizzi IP \\
  HTTP / \textcolor{green-gradient-6}{\textbf{HTTPS}} & \textit{HyperText Transfer Protocol (Secure)}: trasferisce le pagine web \\
  SMTP / \textcolor{green-gradient-6}{\textbf{SMTPS}} & \textit{Simple Mail Transfer Protocol (Secure)}: invia i messaggi di posta \\
  SNMP & \textit{Simple Network Management Protocol}: gestisce i dispositivi di rete \\
  Telnet / \textcolor{green-gradient-6}{\textbf{SSH}} & \textit{Teletype Network} / \textit{Secure SHell}: riga di comando su un host remoto \\
  FTP / \textcolor{green-gradient-6}{\textbf{FTPS}} & \textit{File Transfer Protocol (Secure)}: trasferisce file \\
\end{tabular}
\end{center}

% ============================================================
\section{Dalla URL all'indirizzo: il DNS}

Rileggiamo una URL alla luce di quanto visto: ogni sua parte serve a uno dei
passi della comunicazione.

\begin{figure}[H]
  \centering
  \begin{tikzpicture}[
      part/.style={font=\FiraCodeNL\large, inner sep=2pt, outer sep=0pt, rounded corners=2pt},
      lab/.style={font=\Lato\scriptsize, align=center, inner sep=2pt}]
    \node[part, fill=yellow-gradient-1!70!white]       (sch)  {https};
    \node[part, right=0pt of sch]                       (sep)  {://};
    \node[part, fill=green-gradient-1!60!white, right=0pt of sep] (dom) {www.informatica.it};
    \node[part, right=0pt of dom]                       (col)  {:};
    \node[part, fill=accent!18!white, right=0pt of col] (port) {4242};
    \node[part, fill=neutral!14!white, right=0pt of port] (path) {/corsi/tecweb.html};
    \node[lab, above=16pt of sch]  (lsch)  {\textbf{protocollo}\\ quale protocollo applicativo};
    \node[lab, above=16pt of port] (lport) {\textbf{porta}\\ quale \textbf{processo} sul server};
    \draw[wtdash] (lsch.south)  -- (sch.north);
    \draw[wtdash] (lport.south) -- (port.north);
    \node[lab, below=16pt of dom]  (ldom)  {\textbf{nome dell'host}\\ quale macchina (tramite il DNS)};
    \node[lab, below=16pt of path] (lpath) {\textbf{percorso}\\ quale risorsa sul server};
    \draw[wtdash] (ldom.north)  -- (dom.south);
    \draw[wtdash] (lpath.north) -- (path.south);
  \end{tikzpicture}
  \caption{Le parti di una URL e il passo della comunicazione a cui servono.}
\end{figure}

Resta un problema: le persone ricordano \textbf{nomi}, ma la rete instrada i
pacchetti verso \textbf{indirizzi IP}. A tradurre gli uni negli altri è il
\term{DNS} (\textit{Domain Name System}), la rubrica di Internet.

\begin{itemize}
  \item Il DNS \textbf{non} risiede su un unico server centrale: è un database
        \textbf{distribuito} su migliaia di server in tutto il mondo (server
        \textit{root}, server dei TLD, server \textit{autoritativi}), per
        garantire scalabilità e tolleranza ai guasti.
  \item I nomi sono organizzati come un \textbf{albero rovesciato}: dalla radice
        (\code{.}) si scende ai domini di primo livello o TLD (\code{.com},
        \code{.it}), poi ai domini di secondo livello (\code{google.com},
        \code{unina.it}) e ai sottodomini.
  \item Per non sovraccaricare la rete con interrogazioni ripetute, i server DNS
        e i dispositivi locali tengono in \textbf{cache} i risultati recenti per
        un tempo stabilito, il \term{TTL} (\textit{Time To Live}).
\end{itemize}

\begin{figure}[H]
\centering
\begin{tikzpicture}[
    level distance=1.05cm,
    n/.style={wtnodeplain, font=\FiraCodeNL\scriptsize},
    edge from parent/.style={draw=neutral!50!black, line width=0.6pt},
    level 1/.style={sibling distance=3.2cm},
    level 2/.style={sibling distance=2.3cm},
    level 3/.style={sibling distance=2.9cm}]
  \node[n, fill=primary!18!white] {.}
    child { node[n] {com}
      child { node[n] {google.com} }
      child { node[n] {example.com} } }
    child { node[n] {org}
      child { node[n] {mozilla.org} } }
    child { node[n, fill=green-gradient-1!50!white] {it}
      child { node[n, fill=green-gradient-1!50!white] {unina.it}
        child { node[n, fill=green-gradient-1!50!white] {www.unina.it} }
        child { node[n] {docenti.unina.it} } } };
  \node[wtlab, anchor=west, align=left] at (6.6,-0.1) {radice};
  \node[wtlab, anchor=west, align=left] at (6.6,-1.05) {TLD};
  \node[wtlab, anchor=west, align=left] at (6.6,-2.1) {domini di secondo livello};
  \node[wtlab, anchor=west, align=left] at (6.6,-3.15) {sottodomini};
\end{tikzpicture}
\caption{Lo spazio dei nomi DNS: un albero rovesciato. Un nome si legge da
         destra a sinistra risalendo verso la radice (in verde il percorso di
         \code{www.unina.it}).}
\end{figure}

\begin{esempio}{Chi risponde a «www.unina.it?»}
Il browser non interroga direttamente l'intero albero: chiede al proprio
\term{resolver} (di solito quello del provider o della rete locale), che fa il
lavoro per lui scendendo nell'albero un livello alla volta, a meno che la
risposta non sia già in cache.

\begin{center}
\begin{tikzpicture}[
    att/.style={wtnode, minimum width=2.6cm, minimum height=0.85cm, font=\Lato\scriptsize},
    life/.style={wtdash},
    q/.style={-{Stealth[length=5pt,width=4pt]}, line width=0.8pt, draw=primary!70!black},
    r/.style={-{Stealth[length=5pt,width=4pt]}, line width=0.8pt, draw=accent!70!black},
    ml/.style={font=\Lato\scriptsize, midway, above, inner sep=2pt, fill=white},
    num/.style={circle, fill=accent!75!black, text=white, font=\Lato\tiny\bfseries,
      inner sep=1pt, minimum size=11pt}]
  \foreach \x/\n/\t/\f in {0/b/{\textbf{browser}}/accent!8!white,
                           3.6/res/{\textbf{resolver}\\ (con cache)}/primary!8!white,
                           7.2/root/{server \textbf{root}}/primary!18!white,
                           10.8/tld/{server del TLD \code{.it}}/green-gradient-1!50!white,
                           14.4/aut/{server autoritativo\\ di \code{unina.it}}/orange-gradient-1!50!white} {
    \node[att, fill=\f] (\n) at (\x,0) {\t};
    \draw[life] (\n.south) -- (\x,-5.3);
  }
  \draw[q] (0,-1.0) -- node[ml]{\code{www.unina.it}?} (3.6,-1.0);
  \draw[q] (3.6,-1.7) -- node[ml]{\code{www.unina.it}?} (7.2,-1.7);
  \draw[r] (7.2,-2.2) -- node[ml]{chiedi a \code{.it}} (3.6,-2.2);
  \draw[q] (3.6,-2.9) -- node[ml, pos=0.75]{\code{www.unina.it}?} (10.8,-2.9);
  \draw[r] (10.8,-3.4) -- node[ml, pos=0.75]{chiedi a \code{unina.it}} (3.6,-3.4);
  \draw[q] (3.6,-4.1) -- node[ml, pos=0.8]{\code{www.unina.it}?} (14.4,-4.1);
  \draw[r] (14.4,-4.6) -- node[ml, pos=0.8]{ecco l'indirizzo IP} (3.6,-4.6);
  \draw[r] (3.6,-5.1) -- node[ml]{indirizzo IP} (0,-5.1);
  \foreach \y/\n in {-1.0/1, -1.7/2, -2.9/3, -4.1/4, -5.1/5} \node[num] at (-0.45,\y) {\n};
\end{tikzpicture}
\end{center}
Il resolver chiede prima a un server \textbf{root}, che non conosce l'indirizzo
ma sa chi gestisce \code{.it}; il server del TLD \code{.it} sa chi gestisce
\code{unina.it}; il server \textbf{autoritativo} di \code{unina.it} conosce
finalmente l'indirizzo. Il resolver lo restituisce al browser e lo mette in
cache per la durata del TTL: la volta successiva risponderà da solo, in pochi
millisecondi.
\end{esempio}

% ============================================================
\section{HTTP più da vicino}

HTTP stabilisce come i client web chiedono oggetti (ad esempio pagine web) ai
server e come i server li trasferiscono. Quando l'utente chiede un oggetto, ad
esempio facendo clic su un link, il browser invia al server \textbf{richieste}
HTTP per gli oggetti della pagina; il server risponde con \textbf{risposte} HTTP
che li contengono. Il modello è rimasto lo stesso dagli anni Novanta, ma HTTP
continua a \textbf{evolvere} per bilanciare affidabilità e prestazioni.

\begin{figure}[H]
\centering
\begin{tikzpicture}[
    ver/.style={rounded corners=3pt, fill=#1, text=white, font=\Lato\scriptsize\bfseries,
      minimum width=1.6cm, minimum height=0.6cm},
    des/.style={font=\Lato\small, anchor=west},
    anno/.style={font=\Lato\scriptsize, text=neutral!40!black, anchor=east}]
  \foreach \i/\v/\col/\a/\d in {
      0/{HTTP/1.0}/neutral!55!black/1996/{connessioni spesso \textbf{non persistenti}},
      1/{HTTP/1.1}/accent!75!black/1997/{connessioni \textbf{persistenti} per impostazione predefinita},
      2/{HTTP/2}/orange-gradient-6/2015/{più \textbf{flussi multiplexati} su un'unica connessione TCP},
      3/{HTTP/3}/green-gradient-6/2022/{HTTP su \textbf{QUIC}, che a sua volta usa \textbf{UDP}}} {
    \node[anno] at (-0.2,-\i*0.85) {\a};
    \node[ver=\col, anchor=west] at (0,-\i*0.85) {\v};
    \node[des] at (1.85,-\i*0.85) {\d};
    \draw[neutral!25!white, line width=0.6pt] (1.85,-\i*0.85-0.35) -- (11.5,-\i*0.85-0.35);
  }
\end{tikzpicture}
\caption{Le versioni di HTTP. Richieste e risposte restano il cuore del
         protocollo; cambiano il trasporto e il modo di condividere le
         connessioni.}
\end{figure}

\subsection{TCP o UDP}

HTTP si appoggia principalmente su \textbf{TCP}, ma dal 2022, con HTTP/3, può
usare anche \textbf{UDP}. I due protocolli di trasporto offrono servizi molto
diversi.

\begin{center}
\renewcommand{\arraystretch}{1.4}
\begin{tabular}{@{}>{\raggedright\arraybackslash}p{3.4cm}>{\raggedright\arraybackslash}p{5.6cm}>{\raggedright\arraybackslash}p{5.6cm}@{}}
  & \textbf{TCP} (\textit{Transmission Control Protocol}) & \textbf{UDP} (\textit{User Datagram Protocol}) \\
  \textbf{Connessione} & \textbf{orientato alla connessione}: prima di scambiare dati, i
    processi fanno un \textit{handshake} & \textbf{senza connessione}: nessun handshake, si invia e basta \\
  \textbf{Affidabilità} & trasferimento \textbf{affidabile}: i dati arrivano integri e in
    ordine (se possibile) & \textbf{nessuna garanzia} che i messaggi arrivino, né in che ordine \\
  \textbf{Costo} & più lento: handshake, conferme, ritrasmissioni & leggero, servizi minimi \\
\end{tabular}
\end{center}

Su TCP, il client apre prima una connessione con il server; una volta stabilita,
i processi del browser e del server comunicano attraverso i loro socket. Il
client scrive le richieste HTTP nel proprio socket e ne legge le risposte; il
server legge le richieste dal suo socket e vi scrive le risposte.

HTTP ha usato storicamente TCP proprio per l'\textbf{affidabilità}: TCP
garantisce che richieste e risposte arrivino, prima o poi, integre a
destinazione. Ma i servizi di TCP rallentano la comunicazione, e per questo le
implementazioni recenti di HTTP usano anche UDP, insieme al protocollo
\term{QUIC} (\textit{Quick UDP Internet Connection}), che implementa
l'affidabilità \textbf{sopra} UDP. Le stime riportate a lezione parlano di
comunicazioni circa il 30\% più veloci rispetto a TCP, ma il guadagno reale
dipende molto dalle condizioni della rete: è maggiore sulle reti lente o che
perdono pacchetti.

\info{Il vantaggio di un'architettura a livelli}{
  Qui si vede uno dei grandi vantaggi della stratificazione: le applicazioni
  basate su HTTP non devono preoccuparsi di raggiungibilità, perdita di dati,
  ritrasmissioni. Delegano quasi tutto il lavoro ai livelli inferiori, e restano
  più semplici sia logicamente sia computazionalmente. Passare da TCP a QUIC
  cambia il trasporto, ma non il modo in cui scriviamo una richiesta HTTP.}

\subsection{QUIC}

QUIC porta tre miglioramenti principali.

\textbf{1. Meno attese all'apertura.} Con TCP e TLS, l'apertura della
connessione e la negoziazione della sicurezza avvengono in \textbf{due fasi
separate}, ciascuna con il suo scambio di messaggi. QUIC le combina: di solito
il client può inviare la richiesta dopo un solo \textit{round trip} (andata e
ritorno), e riprendendo una connessione precedente può addirittura inviarla
subito.

\begin{figure}[H]
\centering
\begin{tikzpicture}[
    att/.style={wtnodeplain, minimum width=1.5cm, font=\Lato\scriptsize\bfseries},
    vita/.style={wtdash},
    tcp/.style={-{Stealth[length=4pt,width=3.5pt]}, line width=0.7pt, draw=primary!70!black},
    tls/.style={-{Stealth[length=4pt,width=3.5pt]}, line width=0.7pt, draw=orange-gradient-6},
    req/.style={-{Stealth[length=4pt,width=3.5pt]}, line width=1.1pt, draw=green-gradient-6},
    ml/.style={font=\FiraCodeNL\tiny, sloped, above, inner sep=1pt, fill=white}]
  % --- TCP + TLS ---
  \node[font=\Lato\small\bfseries] at (2.0,0.75) {TCP + TLS 1.3};
  \node[att] (c) at (0,0) {client}; \node[att] (s) at (4,0) {server};
  \draw[vita] (c.south) -- (0,-5.2); \draw[vita] (s.south) -- (4,-5.2);
  \draw[tcp] (0,-0.6) -- node[ml]{SYN} (4,-1.1);
  \draw[tcp] (4,-1.1) -- node[ml]{SYN-ACK} (0,-1.6);
  \draw[tls] (0,-1.7) -- node[ml]{ACK + ClientHello} (4,-2.2);
  \draw[tls] (4,-2.2) -- node[ml]{ServerHello\ldots} (0,-2.7);
  \draw[req] (0,-2.8) -- node[ml]{richiesta HTTP} (4,-3.3);
  \draw[wtbrace] (-0.3,-0.6) -- (-0.3,-2.7) node[midway, left=3pt, font=\Lato\scriptsize, align=right]{2 round trip\\ prima della\\ richiesta};
  % --- QUIC ---
  \begin{scope}[shift={(6.6,0)}]
    \node[font=\Lato\small\bfseries] at (2.0,0.75) {QUIC};
    \node[att] (c2) at (0,0) {client}; \node[att] (s2) at (4,0) {server};
    \draw[vita] (c2.south) -- (0,-5.2); \draw[vita] (s2.south) -- (4,-5.2);
    \draw[tls] (0,-0.6) -- node[ml]{Initial (trasporto + TLS)} (4,-1.1);
    \draw[tls] (4,-1.1) -- node[ml]{risposta (trasporto + TLS)} (0,-1.6);
    \draw[req] (0,-1.7) -- node[ml]{richiesta HTTP} (4,-2.2);
    \draw[wtbrace] (-0.3,-0.6) -- (-0.3,-1.6) node[midway, left=3pt, font=\Lato\scriptsize, align=right]{1 round\\ trip};
    \draw[req, dashed] (0,-3.4) -- node[ml]{0-RTT: richiesta subito} (4,-3.9);
    \node[font=\Lato\tiny, text=neutral!40!black, anchor=north] at (2,-4.2) {riprendendo una connessione precedente};
  \end{scope}
\end{tikzpicture}
\caption{Quanto si aspetta prima di poter inviare la richiesta. In blu
         l'apertura della connessione TCP, in arancione la negoziazione
         crittografica, in verde la richiesta.}
\end{figure}

\textbf{2. Nessun blocco fra richieste indipendenti.} HTTP/2 invia molte
richieste in parallelo sulla \textbf{stessa} connessione TCP. Ma TCP vede un
unico flusso di byte da consegnare in ordine: se si perde un pacchetto, tutto
ciò che arriva dopo resta fermo finché quel pacchetto non viene ritrasmesso,
anche se appartiene a richieste che non c'entrano nulla. È il cosiddetto
\term{head-of-line blocking}. HTTP/3 usa \textbf{flussi QUIC separati}: i dati
mancanti di un flusso non bloccano la consegna degli altri.

\begin{figure}[H]
\centering
\begin{tikzpicture}[
    pk/.style={rectangle, draw=#1!70!black, fill=#1!35!white, line width=0.5pt,
      minimum width=0.55cm, minimum height=0.42cm, font=\FiraCodeNL\tiny},
    persa/.style={rectangle, draw=red-gradient-6, dashed, line width=0.8pt,
      minimum width=0.55cm, minimum height=0.42cm, font=\FiraCodeNL\tiny, text=red-gradient-6},
    ferma/.style={rectangle, draw=neutral!50!white, fill=neutral!12!white, line width=0.5pt,
      minimum width=0.55cm, minimum height=0.42cm, font=\FiraCodeNL\tiny, text=neutral!50!black}]
  % HTTP/2 su TCP
  \node[font=\Lato\small\bfseries, anchor=west] at (0,0.75) {HTTP/2 su TCP: un solo tubo};
  \draw[rounded corners=4pt, draw=neutral!45!black, line width=0.7pt] (0,-0.35) rectangle (5.0,0.35);
  \node[pk=primary] at (0.5,0) {A1};
  \node[persa] at (1.15,0) {A2};
  \foreach \x/\l/\c in {1.8/B1/accent, 2.45/C1/green-gradient-6, 3.1/A3/primary, 3.75/B2/accent, 4.4/C2/green-gradient-6}
    \node[ferma] at (\x,0) {\l};
  \node[font=\Lato\scriptsize, text=red-gradient-6, anchor=north] at (3.45,-0.45) {tutti fermi finché A2 non arriva};
  % HTTP/3 su QUIC
  \begin{scope}[shift={(7.8,0)}]
    \node[font=\Lato\small\bfseries, anchor=west] at (0,0.75) {HTTP/3 su QUIC: flussi separati};
    \foreach \y/\c/\a/\b/\s in {0.25/primary/A1/A2/fermo, -0.25/accent/B1/B2/ok, -0.75/green-gradient-6/C1/C2/ok} {
      \draw[rounded corners=3pt, draw=neutral!45!black, line width=0.6pt] (0,\y-0.22) rectangle (3.4,\y+0.22);
    }
    \node[pk=primary, minimum height=0.36cm] at (0.45,0.25) {A1};
    \node[persa, minimum height=0.36cm] at (1.1,0.25) {A2};
    \node[ferma, minimum height=0.36cm] at (1.75,0.25) {A3};
    \node[pk=accent, minimum height=0.36cm] at (0.45,-0.25) {B1};
    \node[pk=accent, minimum height=0.36cm] at (1.1,-0.25) {B2};
    \node[pk=green-gradient-6, minimum height=0.36cm] at (0.45,-0.75) {C1};
    \node[pk=green-gradient-6, minimum height=0.36cm] at (1.1,-0.75) {C2};
    \node[font=\Lato\scriptsize, anchor=west, text=red-gradient-6] at (3.5,0.25) {aspetta solo A};
    \node[font=\Lato\scriptsize, anchor=west, text=green-gradient-6] at (3.5,-0.25) {consegnato};
    \node[font=\Lato\scriptsize, anchor=west, text=green-gradient-6] at (3.5,-0.75) {consegnato};
  \end{scope}
\end{tikzpicture}
\caption{Tre richieste (A, B, C) in parallelo e un pacchetto perso (A2). Su TCP
         la perdita blocca tutto; su QUIC blocca solo il flusso a cui
         appartiene.}
\end{figure}

\textbf{3. Più facile da aggiornare.} Come TCP, QUIC usa \textbf{conferme}
(\textit{acknowledgment}) e \textbf{timer} per accorgersi delle perdite e
ritrasmettere i dati mancanti. Ma TCP è implementato nel kernel del sistema
operativo, mentre QUIC di solito è implementato nel \textbf{software
applicativo} (il browser, il server): i suoi meccanismi di recupero si possono
aggiornare molto più facilmente.

\subsection{Connessioni persistenti e non persistenti}

In molte applicazioni, HTTP compreso, client e server comunicano a lungo,
scambiandosi molte coppie richiesta--risposta: una dopo l'altra, a intervalli
regolari o in modo intermittente. Con un protocollo orientato alla connessione
(TCP o QUIC) ci sono due possibilità.

\begin{itemize}
  \item \textbf{Connessioni persistenti}: tutte le richieste e le relative
        risposte viaggiano sulla \textbf{stessa} connessione.
  \item \textbf{Connessioni non persistenti}: per ogni coppia
        richiesta--risposta si apre una \textbf{nuova} connessione.
\end{itemize}

Oggi le applicazioni HTTP usano per impostazione predefinita connessioni
persistenti (si possono configurare client e server per non farlo, se serve);
le prime versioni di HTTP (HTTP/1.0) usavano invece connessioni non
persistenti. È ciò che chiede l'header \code{Connection: keep-alive} visto nel
capitolo 1.

\begin{figure}[H]
\centering
\begin{tikzpicture}[x=0.95cm,
    con/.style={draw=primary!60!black, fill=primary!25!white, line width=0.5pt, minimum height=0.42cm},
    req/.style={draw=accent!60!black, fill=accent!25!white, line width=0.5pt, minimum height=0.42cm},
    lab/.style={font=\Lato\scriptsize, anchor=east}]
  \node[lab] at (-0.2,0.6) {non persistente};
  \foreach \i/\o in {0/pagina.html, 1/foto.jpg, 2/stile.css} {
    \draw[con] (\i*3.3,0.39) rectangle (\i*3.3+1.2,0.81);
    \draw[req] (\i*3.3+1.2,0.39) rectangle (\i*3.3+2.9,0.81);
    \node[font=\FiraCodeNL\tiny] at (\i*3.3+2.05,0.6) {\o};
  }
  \node[lab] at (-0.2,-0.4) {persistente};
  \draw[con] (0,-0.61) rectangle (1.2,-0.19);
  \foreach \i/\o in {0/pagina.html, 1/foto.jpg, 2/stile.css} {
    \draw[req] (1.2+\i*1.7,-0.61) rectangle (2.9+\i*1.7,-0.19);
    \node[font=\FiraCodeNL\tiny] at (2.05+\i*1.7,-0.4) {\o};
  }
  \draw[-{Stealth[length=5pt,width=4pt]}, line width=0.7pt, draw=neutral!45!black] (0,-1.0) -- (10.2,-1.0)
    node[font=\Lato\scriptsize, right]{t};
  \node[font=\Lato\scriptsize, anchor=west] at (0,-1.4) {\textcolor{primary!70!black}{\rule{7pt}{7pt}} apertura della connessione
    \qquad \textcolor{accent!70!black}{\rule{7pt}{7pt}} richiesta e risposta};
  \draw[wtbrace] (6.3,-0.65) -- (6.3,-0.15);
  \node[font=\Lato\scriptsize, text=green-gradient-6, anchor=west] at (6.6,-0.4) {si risparmiano due aperture};
\end{tikzpicture}
\caption{Tre oggetti della stessa pagina con connessioni non persistenti e
         persistenti.}
\end{figure}

% ============================================================
\section{I messaggi HTTP in dettaglio}

HTTP ha due formati di messaggio: la \textbf{richiesta}, che specifica il
comando (il metodo) che il server deve eseguire, e la \textbf{risposta}, che
riporta l'esito del comando (successo, fallimento\ldots) ed eventualmente i dati
(ad esempio il file richiesto). In HTTP/1.x entrambi i messaggi sono scritti in
semplice \textbf{testo ASCII}, facilmente leggibile anche da una persona.

\info{E in HTTP/2 e HTTP/3?}{
  Da HTTP/2 in poi i messaggi viaggiano in un formato \textbf{binario}, diviso in
  \textit{frame}, che è più compatto e permette il multiplexing. Il
  \textbf{significato} però non cambia: metodo, destinazione, header, corpo,
  codice di stato sono gli stessi. Per questo i DevTools e strumenti come
  \code{curl} li mostrano sempre nella forma testuale che vediamo qui.}

\subsection{La richiesta}

\begin{figure}[H]
\centering
\begin{tikzpicture}[
    cod/.style={font=\FiraCodeNL\footnotesize, anchor=west, text=white, inner sep=1pt},
    ann/.style={rounded corners=3pt, draw=#1, line width=0.8pt, fill=white, text width=4.6cm,
      align=left, font=\Lato\scriptsize, inner sep=4pt, anchor=west},
    tit/.style={font=\Lato\scriptsize\bfseries, text=#1},
    col/.style={-{Stealth[length=4pt,width=3.5pt]}, line width=0.6pt, draw=#1}]
  \fill[rounded corners=4pt, fill=black!82] (0,0.45) rectangle (8.9,-3.6);
  \node[cod, text=accent!25!white] (l1) at (0.25,0) {GET /courses/techweb?year=2026 HTTP/1.1};
  \node[cod] (l2) at (0.25,-0.5) {Host: www.example.org};
  \node[cod] at (0.25,-0.9) {Accept: text/html};
  \node[cod] at (0.25,-1.3) {Accept-Language: en};
  \node[cod] (l5) at (0.25,-1.7) {User-Agent: Mozilla/5.0};
  \draw[dashed, draw=neutral!55!white, line width=0.6pt] (0.25,-2.15) -- (3.0,-2.15);
  \node[cod, text=neutral!65!white] (l7) at (0.25,-2.75) {[corpo facoltativo, a seconda del metodo]};
  \node[ann=accent!70!black] (a1) at (9.6,0.1) {{\Lato\scriptsize\bfseries\color{accent!70!black}RIGA INIZIALE}\\ metodo + destinazione + versione};
  \node[ann=orange-gradient-6] (a2) at (9.6,-1.1) {{\Lato\scriptsize\bfseries\color{orange-gradient-6}HEADER}\\ metadati e preferenze del client};
  \node[ann=green-gradient-6] (a3) at (9.6,-2.15) {{\Lato\scriptsize\bfseries\color{green-gradient-6}RIGA VUOTA}\\ separa gli header dal corpo};
  \node[ann=primary!70!black] (a4) at (9.6,-3.1) {{\Lato\scriptsize\bfseries\color{primary!70!black}CORPO}\\ i dati della richiesta};
  \draw[col=accent!70!black] (a1.west) -- (8.7,0);
  \draw[col=orange-gradient-6] (a2.west) -- (8.7,-1.1);
  \draw[col=green-gradient-6] (a3.west) -- (8.7,-2.15);
  \draw[col=primary!70!black] (a4.west) -- (8.7,-2.75);
\end{tikzpicture}
\caption{La struttura di una richiesta HTTP (la riga vuota è resa con una
         linea tratteggiata).}
\end{figure}

Il punto interrogativo separa il percorso dalla \textbf{query string}
(\code{year=2026}), che porta parametri aggiuntivi. Il metodo dice al server
\textbf{che cosa fare} con la risorsa indicata.

\subsection{Metodi sicuri e idempotenti}

Prima di passare in rassegna i metodi servono due proprietà, già accennate nel
capitolo 1.

\dfn{Metodo sicuro e metodo idempotente}{
  Un metodo è \term{sicuro} (\textit{safe}) se è di \textbf{sola lettura}: il
  client non chiede alcuna modifica sul server.

  Un metodo è \term{idempotente} se l'\textbf{effetto} di più richieste identiche
  è lo stesso di una sola. Conta lo stato del server, non la risposta: la seconda
  \code{DELETE} della stessa risorsa riceverà probabilmente \code{404}, ma la
  risorsa resta cancellata esattamente come dopo la prima.}

Le due proprietà non sono formalità: sono la ragione per cui un browser, un
\textit{proxy} o una libreria possono \textbf{ripetere} da soli una richiesta
idempotente interrotta da un problema di rete, o \textbf{mettere in cache} la
risposta a una richiesta sicura. Una \code{POST} ripetuta per errore, invece,
può creare due ordini.

\begin{center}
\renewcommand{\arraystretch}{1.35}
\begin{tabular}{@{}l>{\raggedright\arraybackslash}p{6.4cm}ccl@{}}
  \textbf{Metodo} & \textbf{A che cosa serve} & \textbf{Sicuro} & \textbf{Idempotente} & \textbf{Risultato tipico} \\
  \code{GET}     & ottenere una rappresentazione della risorsa         & \mok  & \mok  & \code{200} \\
  \code{HEAD}    & come \code{GET}, ma senza corpo nella risposta      & \mok  & \mok  & \code{200}, \code{304} \\
  \code{POST}    & inviare dati o avviare un'operazione                & \merr & \merr & \code{201} \\
  \code{PUT}     & creare o sostituire la risorsa                       & \merr & \mok  & \code{200}, \code{201}, \code{204} \\
  \code{PATCH}   & applicare una modifica parziale                      & \merr & dipende & \code{200}, \code{204} \\
  \code{DELETE}  & chiedere la rimozione della risorsa                  & \merr & \mok  & \code{204} \\
  \code{OPTIONS} & scoprire le opzioni disponibili per la risorsa       & \mok  & \mok  & \code{204} \\
  \code{QUERY}   & interrogazione sicura con contenuto nel corpo        & \mok  & \mok  & \code{200} \\
\end{tabular}
\end{center}

\subsection{I metodi, uno per uno}

\begin{affianco}[0.55][c]
\begin{lstlisting}[style=http, name=GET]
GET /api/courses/techweb/lectures?week=4 HTTP/1.1
Host: www.example.org
Accept: application/json
Authorization: Bearer <token>
\end{lstlisting}
\risultato
\wtmetodo{GET}{accent!75!black}{Ottiene una rappresentazione della risorsa
  indicata.}{Sicuro e idempotente. Di norma la richiesta non ha corpo.}{200 OK}
\end{affianco}

\begin{affianco}[0.55][c]
\begin{lstlisting}[style=http, name=HEAD]
HEAD /assets/lecture-04.pdf HTTP/1.1
Host: www.example.org
Accept: application/pdf
If-None-Match: "abc123"
\end{lstlisting}
\risultato
\wtmetodo{HEAD}{accent!75!black}{Chiede gli stessi header di una \code{GET},
  ma senza il corpo della risposta.}{Sicuro e idempotente. Utile per leggere i
  metadati (dimensione, data) e per i controlli di cache.}{200 OK
  oppure 304 Not Modified}
\end{affianco}

\info{{If-None-Match, ETag e 304}}{
  Il server può associare a ogni versione di una risorsa un'etichetta, l'header
  \code{ETag} (qui \code{"abc123"}). Il client la ricorda e la rimanda con
  \code{If-None-Match}: se la risorsa non è cambiata, il server risponde
  \code{304 Not Modified} senza ritrasmetterla, e il client usa la copia che ha
  già in cache.}

\begin{affianco}[0.55][c]
\begin{lstlisting}[style=http, name=POST]
POST /api/assignments HTTP/1.1
Host: www.example.org
Content-Type: application/json
Authorization: Bearer <token>

{"title":"HTTP lab","dueDate":"2026-11-05"}
\end{lstlisting}
\risultato
\wtmetodo{POST}{orange-gradient-6}{Invia dati da elaborare; spesso crea una
  nuova risorsa subordinata (qui un nuovo compito).}{Potenzialmente non
  sicuro e, per impostazione predefinita, non idempotente.}{201 Created}
\end{affianco}

\begin{affianco}[0.55][c]
\begin{lstlisting}[style=http, name=PUT]
PUT /api/assignments/42 HTTP/1.1
Host: www.example.org
Content-Type: application/json
Authorization: Bearer <token>

{"id":42,"title":"HTTP lab","dueDate":"2026-11-05"}
\end{lstlisting}
\risultato
\wtmetodo{PUT}{orange-gradient-6}{Crea o sostituisce lo stato della risorsa
  all'URI indicato.}{Non sicuro ma idempotente. Il corpo contiene il
  \textbf{nuovo stato completo}.}{200 OK, 201 Created o 204 No Content}
\end{affianco}

\begin{affianco}[0.55][c]
\begin{lstlisting}[style=http, name=PATCH]
PATCH /api/assignments/42 HTTP/1.1
Host: www.example.org
Content-Type: application/merge-patch+json
Authorization: Bearer <token>

{"dueDate":"2026-11-12"}
\end{lstlisting}
\risultato
\wtmetodo{PATCH}{orange-gradient-6}{Applica una modifica \textbf{parziale}
  allo stato attuale della risorsa.}{Non sicuro. L'idempotenza dipende dal
  documento di patch e da come lo interpreta il server.}{200 OK o 204 No
  Content}
\end{affianco}

\begin{affianco}[0.55][c]
\begin{lstlisting}[style=http, name=DELETE]
DELETE /api/assignments/42 HTTP/1.1
Host: www.example.org
Authorization: Bearer <token>
If-Match: "assignment-v3"
\end{lstlisting}
\risultato
\wtmetodo{DELETE}{red-gradient-6}{Chiede di rimuovere l'associazione fra
  l'URI e la sua risorsa.}{Non sicuro ma idempotente: ripeterla ha lo stesso
  effetto voluto.}{204 No Content}
\end{affianco}

\begin{esempio}{POST, PUT e PATCH ripetute}
Che cosa succede se la stessa richiesta parte \textbf{due volte}, per esempio
perché l'utente ha fatto doppio clic o la rete ha ritrasmesso?

\begin{center}
\renewcommand{\arraystretch}{1.4}
\begin{tabular}{@{}>{\raggedright\arraybackslash}p{6.4cm}>{\raggedright\arraybackslash}p{4.6cm}>{\raggedright\arraybackslash}p{5.4cm}@{}}
  \textbf{Richiesta} & \textbf{Dopo la prima} & \textbf{Dopo la seconda} \\
  \leavevmode\code{POST /api/assignments} & esiste il compito 43 & esistono i compiti 43
    \textbf{e 44}: un duplicato \\
  \leavevmode\code{PUT /api/assignments/42} con lo stato completo & il compito 42 ha quello
    stato & identico: \textbf{nessun cambiamento} \\
  \leavevmode\code{PATCH} con \code{\{"dueDate":"2026-11-12"\}} \newline (imposta un valore) & la scadenza è il 12 &
    identica: questa patch è idempotente \\
  \leavevmode\code{PATCH} che \textbf{aggiunge} un elemento a una lista & un elemento in
    più & \textbf{due} elementi in più: non idempotente \\
  \leavevmode\code{DELETE /api/assignments/42} & il compito non esiste più (\code{204}) &
    ancora inesistente, anche se la risposta è \code{404} \\
\end{tabular}
\end{center}
È per questo che \code{PATCH} è indicato come «dipende»: una patch che
\textit{imposta} dei valori (come il \textit{JSON Merge Patch} dell'esempio) è
idempotente, una che \textit{aggiunge} o \textit{incrementa} no.
\end{esempio}

\begin{affianco}[0.55][c]
\begin{lstlisting}[style=http, name=OPTIONS]
OPTIONS /api/assignments HTTP/1.1
Host: www.example.org
Origin: https://app.example.org
Access-Control-Request-Method: POST
Access-Control-Request-Headers: content-type
\end{lstlisting}
\risultato
\wtmetodo{OPTIONS}{neutral!55!black}{Chiede quali opzioni di comunicazione
  sono disponibili per la risorsa.}{Sicuro e idempotente. I browser lo usano
  anche per il \textit{preflight} CORS.}{204 No Content, con gli header
  \code{Allow} o CORS}
\end{affianco}

L'esempio è proprio un \textit{preflight}: prima di inviare una \code{POST}
con corpo JSON da \code{app.example.org} a \code{www.example.org}, il browser
chiede al server se l'operazione è permessa da quell'origine. Lo vedremo nel
capitolo 23, parlando di \textit{same-origin policy} e CORS.

\begin{affianco}[0.55][c]
\begin{lstlisting}[style=http, name=QUERY]
QUERY /api/assignments/search HTTP/1.1
Host: www.example.org
Content-Type: application/json
Accept: application/json

{"course":"TechWeb","status":"open"}
\end{lstlisting}
\risultato
\wtmetodo{QUERY}{green-gradient-6}{Esegue sul server un'interrogazione
  descritta dal contenuto della richiesta.}{Sicuro, idempotente e
  memorizzabile in cache. L'header \code{Content-Type} è obbligatorio.}{200 OK}
\end{affianco}

\info{Perché serviva un metodo nuovo}{
  \code{QUERY} è il metodo più recente (le slide citano la RFC 10008) e
  colma un vuoto. Una ricerca complessa con \code{GET} deve stare tutta nella
  query string, che ha limiti di lunghezza e finisce nei log; la si può mettere
  nel corpo di una \code{POST}, ma allora la richiesta non è più né sicura né
  idempotente, e non si può mettere in cache. \code{QUERY} unisce i due mondi: un
  corpo come la \code{POST}, la semantica di sola lettura della \code{GET}.
  Essendo così recente, il supporto da parte di server e librerie può essere
  ancora limitato.}

\subsection{La risposta}

\begin{figure}[H]
\centering
\begin{tikzpicture}[
    cod/.style={font=\FiraCodeNL\footnotesize, anchor=west, text=white, inner sep=1pt},
    ann/.style={rounded corners=3pt, draw=#1, line width=0.8pt, fill=white, text width=4.6cm,
      align=left, font=\Lato\scriptsize, inner sep=4pt, anchor=west},
    tit/.style={font=\Lato\scriptsize\bfseries, text=#1},
    col/.style={-{Stealth[length=4pt,width=3.5pt]}, line width=0.6pt, draw=#1}]
  \fill[rounded corners=4pt, fill=black!82] (0,0.45) rectangle (8.9,-3.6);
  \node[cod, text=green-gradient-1] at (0.25,0) {HTTP/1.1 200 OK};
  \node[cod] at (0.25,-0.5) {Date: Mon, 12 Oct 2026 10:15:00 GMT};
  \node[cod] at (0.25,-0.9) {Content-Type: application/json};
  \node[cod] at (0.25,-1.3) {Content-Length: 38};
  \node[cod] at (0.25,-1.7) {Cache-Control: max-age=60};
  \draw[dashed, draw=neutral!55!white, line width=0.6pt] (0.25,-2.15) -- (3.0,-2.15);
  \node[cod, text=orange-gradient-1] at (0.25,-2.75) {\{"course":"TechWeb","status":"online"\}};
  \node[ann=green-gradient-6] (a1) at (9.6,0.1) {{\Lato\scriptsize\bfseries\color{green-gradient-6}RIGA DI STATO}\\ versione + codice + messaggio};
  \node[ann=orange-gradient-6] (a2) at (9.6,-1.1) {{\Lato\scriptsize\bfseries\color{orange-gradient-6}HEADER}\\ metadati della risposta};
  \node[ann=accent!70!black] (a3) at (9.6,-2.15) {{\Lato\scriptsize\bfseries\color{accent!70!black}RIGA VUOTA}\\ chiude la sezione degli header};
  \node[ann=primary!70!black] (a4) at (9.6,-3.1) {{\Lato\scriptsize\bfseries\color{primary!70!black}CORPO}\\ la rappresentazione o il risultato};
  \draw[col=green-gradient-6] (a1.west) -- (8.7,0);
  \draw[col=orange-gradient-6] (a2.west) -- (8.7,-1.1);
  \draw[col=accent!70!black] (a3.west) -- (8.7,-2.15);
  \draw[col=primary!70!black] (a4.west) -- (8.7,-2.75);
\end{tikzpicture}
\caption{La struttura di una risposta HTTP.}
\end{figure}

\code{Content-Length} vale \code{38} perché il corpo è lungo 38 byte, e
\code{Cache-Control: max-age=60} autorizza il client a riusare la risposta per
60 secondi senza richiederla.

\begin{center}
\begin{tikzpicture}[
    cls/.style={rounded corners=3pt, fill=#1, text=white, font=\Lato\scriptsize\bfseries,
      minimum width=1.2cm, minimum height=0.55cm},
    nome/.style={font=\Lato\small\bfseries, anchor=west},
    ex/.style={font=\FiraCodeNL\small, anchor=west}]
  \foreach \i/\c/\col/\n/\e in {
      0/1xx/neutral!55!black/Informativi/{l'elaborazione continua},
      1/2xx/green-gradient-6/Successo/{200 OK, 201 Created, 204 No Content},
      2/3xx/primary!75!black/Redirezione/{301, 302, 304},
      3/4xx/orange-gradient-6/Errore del client/{400, 401, 403, 404},
      4/5xx/red-gradient-6/Errore del server/{500, 502, 503}} {
    \node[cls=\col] at (0,-\i*0.75) {\c};
    \node[nome] at (1.0,-\i*0.75) {\n};
    \node[ex, text=\col] at (5.0,-\i*0.75) {\e};
  }
\end{tikzpicture}
\end{center}

Le classi le conosciamo dal capitolo 1. Alcune coppie di codici vicini vengono
confuse spesso, ed è utile saperle distinguere:

\begin{center}
\renewcommand{\arraystretch}{1.35}
\begin{tabular}{@{}l>{\raggedright\arraybackslash}p{12.6cm}@{}}
  \code{301} / \code{302} & la risorsa si è spostata \textbf{per sempre}
    (aggiornare i link e i segnalibri) / \textbf{temporaneamente} (continuare a
    usare il vecchio indirizzo) \\
  \code{304} & non è un errore: «la tua copia in cache è ancora buona» \\
  \code{401} / \code{403} & \textbf{non so chi sei} (manca l'autenticazione, o non
    è valida) / \textbf{so chi sei, ma non puoi} (mancano i permessi) \\
  \code{502} / \code{503} & un server intermedio (\textit{proxy}, \textit{gateway})
    ha ricevuto una risposta non valida da quello dietro / il server è
    \textbf{momentaneamente} non disponibile, ad esempio per sovraccarico o
    manutenzione \\
\end{tabular}
\end{center}

% ============================================================
\section{HTTPS: HTTP dentro un canale TLS}

\term{HTTPS} non è un protocollo diverso da HTTP: è lo \textbf{stesso} HTTP,
trasportato dentro un canale protetto da \term{TLS} (\textit{Transport Layer
Security}), che si inserisce fra HTTP e il trasporto.

\begin{figure}[H]
\centering
\begin{tikzpicture}[
    liv/.style={rounded corners=3pt, fill=#1, text=white, minimum width=10.8cm, minimum height=0.75cm},
    nome/.style={font=\Lato\small\bfseries, text=white, anchor=west},
    des/.style={font=\Lato\small, text=white, anchor=west}]
  \foreach \i/\col/\n/\d in {
      0/accent!75!black/HTTP/{richiesta e risposta},
      1/orange-gradient-6/TLS/{cifra, autentica il server, controlla l'integrità},
      2/green-gradient-6/{TCP o QUIC}/{trasporta i dati},
      3/primary!75!black/IP/{instrada i pacchetti verso la destinazione}} {
    \node[liv=\col, anchor=west] at (0,-\i*0.9) {};
    \node[nome] at (0.3,-\i*0.9) {\n};
    \node[des] at (3.4,-\i*0.9) {\d};
  }
\end{tikzpicture}
\caption{La pila di HTTPS. Con QUIC, TLS è integrato nel protocollo di
         trasporto, ma i servizi offerti sono gli stessi.}
\end{figure}

TLS offre tre garanzie:
\begin{itemize}
  \item \textbf{riservatezza}: i dati sono cifrati, e chi li intercetta lungo il
        percorso vede solo byte incomprensibili;
  \item \textbf{autenticazione}: il server dimostra la propria identità con un
        \term{certificato} firmato da un'autorità di cui il browser si fida;
  \item \textbf{integrità}: se qualcuno modifica i dati in transito, il
        destinatario se ne accorge.
\end{itemize}

\warning{Il lucchetto non vuol dire «sito affidabile»}{
  Il lucchetto nella barra degli indirizzi significa soltanto che il canale è
  \textbf{cifrato} e che il certificato è stato \textbf{verificato} per quel nome
  di dominio. \textbf{Non} dice nulla sull'onestà di chi gestisce il sito:
  ottenere un certificato valido per \code{poste-italiane-verifica.com} è facile
  e gratuito, e moltissimi siti di \textit{phishing} mostrano il lucchetto come
  tutti gli altri.}

% ============================================================
\section{Dal clic alla pagina: il viaggio completo}

Mettiamo insieme tutti i pezzi. Quando facciamo clic su un link
\code{https://}, prima che compaia la pagina il browser compie sei passi.

\begin{figure}[H]
\centering
\begin{tikzpicture}[
    passo/.style={rounded corners=3pt, draw=#1, line width=0.9pt, fill=white,
      minimum width=4.7cm, minimum height=1.25cm, text width=4.3cm, align=left, inner sep=6pt},
    num/.style={font=\Lato\Large\bfseries, text=#1},
    tit/.style={font=\Lato\small\bfseries},
    des/.style={font=\Lato\scriptsize, text=neutral!35!black},
    fr/.style={-{Stealth[length=5pt,width=4pt]}, line width=0.8pt, draw=neutral!50!black}]
  \foreach \i/\x/\y/\col/\t/\d in {
      1/0/0/accent!75!black/{ANALISI DELLA URL}/{schema, host, porta, percorso},
      2/5.2/0/green-gradient-6/{DNS}/{dal nome all'indirizzo IP},
      3/10.4/0/orange-gradient-6/{CONNESSIONE}/{TCP o QUIC verso IP e porta},
      4/0/-1.75/primary!75!black/{TLS}/{certificato e chiavi di sessione},
      5/5.2/-1.75/accent!75!black/{RICHIESTA}/{metodo, destinazione, header, corpo},
      6/10.4/-1.75/green-gradient-6/{RISPOSTA}/{stato, header, corpo}} {
    \node[passo=\col] (p\i) at (\x,\y)
      {{\Lato\Large\bfseries\color{\col}\i}\quad{\Lato\small\bfseries\t}\\[2pt]
       {\Lato\scriptsize\color{neutral!35!black}\d}};
  }
  \draw[fr] (p1) -- (p2); \draw[fr] (p2) -- (p3);
  \draw[fr] (p3.south) -- ++(0,-0.25) -| (p4.north);
  \draw[fr] (p4) -- (p5); \draw[fr] (p5) -- (p6);
\end{tikzpicture}
\caption{Il viaggio completo, dal clic alla pagina. Il browser fa tutto questo
         mentre noi diciamo «Internet è lento».}
\end{figure}

I primi passi li abbiamo già visti nel capitolo 1; la novità è il passo 4, e il
fatto che ognuno di questi passi si vede, e si misura, nei DevTools. Nella
scheda \textit{Rete}, selezionando una richiesta e aprendo la sezione
\textit{Tempistiche}, il browser mostra quanto è durata ciascuna fase.

\begin{figure}[H]
\centering
\begin{minipage}{15.4cm}
\begin{devtools}{Rete}
  {\Lato\scriptsize
  \noindent\hspace*{6pt}{\color{neutral!30!black}Header\quad Cookie\quad Richiesta\quad Risposta\quad
    \colorbox{dtline}{Tempistiche}\quad Sicurezza}\par\vspace{2pt}
  {\color{dtline}\hrule}\vspace{6pt}
  \noindent\hspace*{8pt}\begin{tikzpicture}[x=0.06cm,
      fase/.style={font=\Lato\scriptsize, anchor=east, text=neutral!20!black},
      ms/.style={font=\FiraCodeNL\scriptsize, anchor=west, text=neutral!30!black}]
    \foreach \i/\n/\a/\d/\col in {
        0/Bloccata/0/2/neutral!40!white,
        1/Risoluzione DNS/2/14/green-gradient-6,
        2/Connessione/16/22/orange-gradient-6,
        3/Configurazione TLS/38/26/primary!70!black,
        4/Invio/64/1/accent!70!black,
        5/In attesa/65/78/accent!45!white,
        6/Ricezione/143/7/primary!45!white} {
      \node[fase] at (-2,-\i*0.42) {\n};
      \fill[\col] (\a,-\i*0.42-0.13) rectangle (\a+\d,-\i*0.42+0.13);
      \node[ms] at (\a+\d+1,-\i*0.42) {\d\ ms};
    }
    \draw[neutral!40!white] (0,0.3) -- (0,-2.8);
  \end{tikzpicture}\par\vspace{6pt}
  {\color{dtline}\hrule}\vspace{2pt}
  \noindent\hspace*{6pt}{\color{neutral!25!black}Totale: 150 ms\quad$\vert$\quad
    valori di esempio}\par\vspace{3pt}}
\end{devtools}
\end{minipage}
\caption{La sezione \textit{Tempistiche} di una richiesta HTTPS (valori di
         esempio). Ogni barra è uno dei passi della figura precedente.}
\end{figure}

Leggerla non è difficile, ora che sappiamo che cosa c'è sotto:
\begin{itemize}
  \item \textbf{Risoluzione DNS}: il passo 2. Alla seconda visita è spesso zero,
        perché l'indirizzo è in cache;
  \item \textbf{Connessione} e \textbf{Configurazione TLS}: i passi 3 e 4, cioè
        i \textit{round trip} visti a proposito di QUIC. Con una connessione
        \textbf{persistente} già aperta spariscono del tutto;
  \item \textbf{Invio}: il tempo per spedire la richiesta, di solito
        trascurabile;
  \item \textbf{In attesa}: il tempo fra l'invio della richiesta e l'arrivo del
        primo byte della risposta (\textit{Time To First Byte}). È il tempo che
        il \textbf{server} ha impiegato a elaborare la richiesta: quando una
        pagina è lenta, il colpevole si trova spesso qui;
  \item \textbf{Ricezione}: il tempo per scaricare il corpo della risposta, che
        dipende dalla sua dimensione e dalla banda.
\end{itemize}

% ============================================================
\section{Riepilogo}

\begin{figure}[H]
\centering
\begin{tikzpicture}[
    idea/.style={rounded corners=4pt, draw=#1, line width=1pt, fill=white, minimum width=3.9cm,
      minimum height=2.0cm, text width=3.6cm, align=center, inner sep=5pt},
    tit/.style={font=\Lato\normalsize\bfseries, text=#1},
    des/.style={font=\Lato\scriptsize, text=neutral!30!black}]
  \foreach \i/\col/\t/\d in {
      0/accent!75!black/{RETE}/{definisce il contesto e la scala: processi, socket, indirizzi IP e porte},
      1/green-gradient-6/{URL + DNS}/{identificano e trovano la destinazione},
      2/orange-gradient-6/{HTTP}/{definisce la conversazione: metodi, header, codici di stato},
      3/primary!75!black/{HTTPS}/{protegge i dati in transito: riservatezza, autenticazione, integrità}} {
    \node[idea=\col] at (\i*4.25,0)
      {{\Lato\normalsize\bfseries\color{\col}\t}\\[4pt]{\Lato\scriptsize\color{neutral!30!black}\d}};
  }
\end{tikzpicture}
\caption{Quattro idee da portare a casa.}
\end{figure}

\begin{itemize}
  \item Un'applicazione di rete è fatta di \textbf{processi} su host diversi, che
        comunicano attraverso i \textbf{socket}; l'indirizzo \textbf{IP} trova
        l'host, la \textbf{porta} trova il processo.
  \item Il \textbf{DNS} traduce i nomi in indirizzi IP con un database
        distribuito e gerarchico, e una cache regolata dal TTL.
  \item HTTP mantiene lo stesso modello richiesta--risposta, ma il trasporto
        evolve: connessioni persistenti (1.1), flussi multiplexati (2), QUIC su
        UDP (3).
  \item I metodi HTTP hanno una semantica precisa: \textbf{sicuri} (\code{GET},
        \code{HEAD}, \code{OPTIONS}, \code{QUERY}) e \textbf{idempotenti}
        (\code{PUT}, \code{DELETE} e tutti i sicuri); \code{POST} non è né l'uno né
        l'altro.
  \item \textbf{HTTPS} è HTTP dentro TLS: dati cifrati, server autenticato,
        integrità controllata. Il lucchetto non garantisce che il sito sia onesto.
\end{itemize}

% ============================================================
\section{Riferimenti}

\begin{itemize}
  \item \textbf{RFC 1034} --- \textit{Domain Names: Concepts and Facilities},
        sezioni 2.3--2.4. \\
        \urlx{https://www.rfc-editor.org/rfc/rfc1034.html}
  \item \textbf{RFC 9110} --- \textit{HTTP Semantics}, sezioni 3, 5--6, 9 e 15. \\
        \urlx{https://www.rfc-editor.org/rfc/rfc9110.html}
  \item \textbf{RFC 9114} --- \textit{HTTP/3}, sezioni 1 e 6.1. \\
        \urlx{https://www.rfc-editor.org/rfc/rfc9114.html}
  \item \textbf{RFC 5789} --- \textit{PATCH Method for HTTP}, sezioni 1--2.2. \\
        \urlx{https://www.rfc-editor.org/rfc/rfc5789.html}
  \item \textbf{RFC 7396} --- \textit{JSON Merge Patch}, sezioni 1--3. \\
        \urlx{https://www.rfc-editor.org/rfc/rfc7396.html}
  \item \textbf{RFC 8446} --- \textit{The Transport Layer Security Protocol Version
        1.3}, sezioni 1--2, 4 e 5. \\
        \urlx{https://www.rfc-editor.org/rfc/rfc8446.html}
  \item \textbf{RFC 10008} --- \textit{The HTTP QUERY Method}, sezioni 1--2 e
        tabella 1. \\
        \urlx{https://www.rfc-editor.org/rfc/rfc10008.html}
  \item \textbf{Fetch Standard} --- WHATWG, sezione \textit{CORS-preflight fetch}. \\
        \urlx{https://fetch.spec.whatwg.org/\#cors-preflight-fetch}
  \item \textbf{RFC 3986} --- \textit{Uniform Resource Identifier: Generic
        Syntax}, sezioni 3--3.5. \\
        \urlx{https://www.rfc-editor.org/rfc/rfc3986.html}
\end{itemize}
